% !TeX program = xelatex
\documentclass[zihao=-4]{ctexart}  % 小四号字

% 中文支持
\usepackage{xeCJK}
\setCJKmainfont{SimSun}[AutoFakeBold, AutoFakeSlant]   % 宋体
\setCJKsansfont{SimHei}[AutoFakeBold]                  % 黑体

% 页面边距
\usepackage[margin=2.5cm]{geometry}

% 必要宏包
\usepackage{graphicx}
\usepackage{ulem}          % 下划线
\usepackage{array}
\usepackage{setspace}

% 取消段落缩进
\setlength{\parindent}{0pt}

% 自定义下划线填空（长度可调）
\newcommand{\fillblank}[1][5cm]{\uline{\hspace{#1}}}
\newcommand{\filldate}{%
    \underline{\makebox[1.5cm]{}}年\underline{\makebox[1cm]{}}月\underline{\makebox[1cm]{}}日%
}

\ctexset{
    section = {
        format = \raggedright\Large\bfseries,   % 左对齐、大号、粗体
        name = {,},                             % 不添加“第”“节”等前缀
        number = \chinese{section},             % 中文数字编号（仅对 \section 有效）
        aftername = \quad,                      % 标题编号与标题文字之间的间距
    }
}

\begin{document}

% ========== 封面 ==========

% 左上角放置两张图片（并排，左对齐）
\noindent
\begin{minipage}{0.2\textwidth}
    \includegraphics[width=1.5cm]{media/image1.png}   % 替换为 media/image1.png
\end{minipage}
\begin{minipage}{0.6\textwidth}
    \includegraphics[width=3cm]{media/image2.png}     % 替换为 media/image2.png
\end{minipage}
\hfill
\begin{minipage}{0.1\textwidth}
    % 留空
\end{minipage}

\vspace{1cm}

% 主标题（居中）
\begin{center}
    \vspace{2cm}
    {\Huge\bfseries 课程名称}\\[0.8cm]
    {\LARGE 实验报告}\\[1.2cm]
    \vspace{1.4cm}
    {\fontsize{18pt}{21.6pt}\selectfont 202X--202X 学年第 X 学期}\\[1.5cm]
    \vspace{1.6cm}
    % 信息表格：使用宋体小四（12pt）
    {\fontsize{14pt}{14.4pt}\selectfont
    \renewcommand{\arraystretch}{1.5}
    \begin{tabular}{p{2.5cm} p{8cm}}
        专业： & \underline{\makebox[7cm][c]{专业名称}} \\[0.3cm]
        班级： & \underline{\makebox[7cm][c]{班级名称}} \\[0.3cm]
        学号： & \underline{\makebox[7cm][c]{12345678}} \\[0.3cm]
        姓名： & \underline{\makebox[7cm][c]{张三}} \\
    \end{tabular}
    }
    
    \vspace{2cm}
    
    {\underline{\makebox[1.5cm]{202X}}年\underline{\makebox[1cm]{X}}月\underline{\makebox[1cm]{X}}日}
    
    \vfill
\end{center}

\newpage

% ========== 正文标题（无编号） ==========
\section*{一、实验题目}
% 实验内容
\begin{center}
CIFAR-10 图像分类深度网络优化和路透社分类改进
\end{center}

\section*{二、实验目的}
% 实验目的
1. 在 CIFAR-10 数据集上构建一个具有 25 个隐藏层、每层 80 个神经元的深度神经网络。

2. 比较不同技术（批归一化、SELU 自归一化、Alpha Dropout + MC Dropout）对模型收敛速度、泛化能力和训练时间的影响。

3. 验证 MC Dropout 能否通过测试时多次随机采样提升分类准确率。

4. 将实验1中的深度网络优化技术应用于路透社新闻数据集（46 个主题的多分类问题）。

\section*{三、实验过程（步骤）}
% 实验步骤
\subsection*{3.1 数据准备}
加载 CIFAR-10 数据集（32×32 彩色图像，10 类），共 60000 张。

图像像素值缩放到 [0,1] 用于 ELU/BN 模型；计算训练集均值和标准差，将输入标准化（均值 0，方差 1）用于 SELU 模型。再将图像展平为 3072 维向量。

从训练集中划分 10\% 作为验证集。

\subsection*{3.2 模型构建}
基础模型：25 个 Dense，输出层 Dense(10, softmax)。使用 Nadam 优化器，学习率 0.001。

批归一化模型：每个全连接层后添加 BatchNormalization，然后接 ELU 激活（无偏置）。

SELU 自归一化模型：全连接层使用 SELU 激活 + LeCun 正态初始化，输入标准化。

Alpha Dropout 模型：在 SELU 模型基础上每层后添加 AlphaDropout(0.1)，训练后使用 MC Dropout 测试（20 次前向传播取平均）。

\subsection*{3.3 训练与评估}
绘制验证准确率与损失曲线对比。

\subsection*{3.4 改进路透社分类算法}

与实验1相同：25 个隐藏层，每层 80 个神经元。

输出层为 46 个神经元的 softmax。

优化器：Nadam，基础模型和 BN 模型学习率 0.001，SELU 和 AlphaDropout 模型学习率 0.0005（因 SELU 对学习率更敏感）。

\section*{四、实验结果与讨论}
% 结果与讨论
\subsection*{4.1 CIFAR-10 图像分类训练图表}
\begin{figure}[htbp]   % 位置参数：here, top, bottom, page
    \centering
    \includegraphics[width=0.8\textwidth]{media/8/Figure_2.png}
    \caption{验证损失随迭代次数的变化曲线}
    \label{fig:loss_curve}
\end{figure}
\begin{figure}[htbp]   % 位置参数：here, top, bottom, page
    \centering
    \includegraphics[width=0.8\textwidth]{media/8/Figure_1.png}
    \caption{验证准确度随迭代次数的变化曲线}
    \label{fig:loss_curve}
\end{figure}

\subsection*{4.2 改进路透社分类训练图表}
\begin{figure}[htbp]   % 位置参数：here, top, bottom, page
    \centering
    \includegraphics[width=0.8\textwidth]{media/8/Figure_4.png}
    \caption{验证损失随迭代次数的变化曲线}
    \label{fig:loss_curve}
\end{figure}
\begin{figure}[htbp]   % 位置参数：here, top, bottom, page
    \centering
    \includegraphics[width=0.8\textwidth]{media/8/Figure_3.png}
    \caption{验证准确度随迭代次数的变化曲线}
    \label{fig:loss_curve}
\end{figure}

\subsection*{4.3 分析}
批归一化是稳定深层网络训练最有效的方法，同时提高了准确率和收敛速度。

SELU 配合标准化输入可作为一种无需 BN 的替代方案，但在极深网络中仍需谨慎。

MC Dropout 无需重新训练即可获得集成效果，提升约 1.2% 的准确率。
\end{document}